\documentclass[12pt]{article}

\usepackage[a4paper, margin=1in, heightrounded]{geometry}
\usepackage{setspace}
\onehalfspace


\title{Reaction Paper: Game Development in 8-bits}
\author{John Jovic A. Gutierrez - 202207198}
\date{February 11, 2026}
\begin{document}
\maketitle
\thispagestyle{empty}

The "Game Development in Eight Bits" talk by Kevin Zurawel was fascinating to where I can relate from my EEE courses such as Computer Architecture and Digital Logic Design (EEE 143/153) as well as Networking and Data Communication (EEE 157) in a way that resources are naturally scarce and limited, hence data must be compressed and optimized to fit within the constraints of the hardware. With this, it was impressive that they employed various clever techniques to achieve features that made the games engaging and enjoyable.

I appreciate how Engineering, in general, must use creativity and innovation to solve such problems given the constraints of the system. With this, appreciation towards my CoE/EEE courses has been heightened as I can see how the concepts and skills I am learning can make further optimized systems even though we are our systems are not as limited as the 8-bit era anymore. Though the era of Big Data and AI made it relevant to be able to optimize our systems once again as computing hardware, specifically RAM or memory are becoming on demand for AI development and processing.

I am the type of person who would like to optimize my systems in a reasonable manner, especially by not letting a large portion of my system's memory get eaten by unnecessary bloat nor processes, which are not in use, hence I can say that the talk was quite insightful in realizing that there are more lower level optimizations that we can do with a more creative approach, such as having more level of abstractions, rather than having to draw everything what is on screen and filling up the memory.

In retrospect, such optimized production took a lot of time and mental fortitude to develop, which adds a barrier for entry in development jobs if it is strict in some way, I can say that higher level development is essential for feature adding and innovation, whereas lower level abstractions optimize the systems higher level development sits on. I also want to appreciate the quote "Embrace the stupid", closely to the saying: "Keep it Simple and Stupid" when developing systems to be more extensible and resource friendly just to get the job done.

\end{document}
